\section{Банк практических заданий}

В этом разделе собраны задачи для закрепления материала. Задания разделены по темам и уровням сложности:
\begin{itemize}
    \item \textbf{Базовый (Б):} Проверка понимания основных команд.
    \item \textbf{Продвинутый (П):} Требует написания алгоритмов.
    \item \textbf{Комплексный (К):} Мини-проект, объединяющий несколько тем.
\end{itemize}

У каждого задания указано, где его выполнять: \textbf{симулятор} (редактор Python на сайте, нижняя камера \verb|Camera(view="bottom")|, модель ArUco), \textbf{компьютер} (Python с полной библиотекой OpenCV) или \textbf{дрон} (настоящий Пионер или веб-камера). Строка «Опора» называет пример из папки \verb|examples/python|, с которого удобно начать. Миссии из строки «В симуляторе» открываются в панели \ui{Обучение}; сцену, которую готовит миссия, можно использовать и для своей программы.

\begin{infobox}[Чего нет в симуляторе]
Браузерный \verb|cv2| --- учебная модель: в нём есть кадр камеры, детектор ArUco, рисование и \verb|imshow|, но нет \verb|cv2.imwrite| (ничего не сохраняет), \verb|solvePnP| (всегда неуспех), цветовых масок, сравнения кадров и распознавания лиц. В \verb|pioneer_sdk| нет углов крена и тангажа и высоты по барометру: есть положение LPS, курс, дальномер, батарея и оптический поток. Задания, которым это нужно, помечены «дрон» или «компьютер».
\end{infobox}

\subsection{Блок 1: Работа с камерой и ArUco}

\begin{taskbox}[(Б) Фотограф]
\emph{Дрон.} Напишите скрипт, который делает фото каждые 5 секунд и сохраняет их в папку, в названии которой --- дата и время запуска (\verb|time.strftime("%Y-%m-%d_%H-%M-%S")|).
\expected{за минуту работы в новой папке 12 файлов JPG с возрастающими номерами; ни один файл не пустой (\texttt{cv2.imread} читает каждый).}
Опора: \verb|camera_examples/take_photo.py|, \verb|camera_examples/get_frames_from_camera.py|.
\end{taskbox}

\begin{taskbox}[(Б) Детектор движения]
\emph{Дрон или веб-камера, компьютер.} Сравнивайте два последовательных серых кадра (\verb|cv2.absdiff|, затем \verb|cv2.threshold|). Если количество изменившихся пикселей превышает порог --- выводите сообщение <<Движение обнаружено!>>.
\expected{на неподвижной сцене сообщений нет 10 секунд подряд; взмах рукой перед камерой даёт сообщение в течение секунды. Порог подобран так, что мерцание ламп его не превышает.}
Опора: \verb|camera_examples/camera_stream.py|.
\end{taskbox}

\begin{taskbox}[(П) Умный светофор]
\emph{Дрон.} Распечатайте красный и зеленый круги. Если дрон видит зеленый цвет --- он взлетает на 1 метр. Если красный --- садится. Цвет ищите маской в пространстве HSV (\verb|cv2.inRange|) и считайте долю пикселей маски.
\expected{зелёный круг перед камерой --- взлёт, красный --- посадка; решение принимается по трём кадрам подряд, поэтому случайный кадр с цветной одеждой не поднимает дрон.}
Опора: \verb|camera_examples/camera_stream.py|, \verb|mission_examples/PiZero_test.py|.
\end{taskbox}

\begin{taskbox}[(П) Поиск маркера]
\emph{Симулятор.} Найдите метку 12 перебором конечного списка точек обзора на высоте 3~м --- квадрат $3 \times 3$ точки с шагом 2~м вокруг старта. В каждой точке снимайте новый кадр; как только метка найдена, прекратите поиск, зависните над ней 3 секунды и вернитесь на старт с посадкой.
\expected{программа печатает номер точки, в которой метка найдена, не посещает точки после находки и при отсутствии метки печатает «Поиск истёк» после девяти точек. На настоящем дроне с передней камерой ту же логику делают вращением по курсу (\texttt{yaw}) с шагом 30°.}
\mission{Python}{85, 101}
Опора: \verb|aruco_examples/11_aruco_basics_example1_detect.py|.
\end{taskbox}

\begin{taskbox}[(П) Дистанция]
\emph{Дрон.} Дрон висит напротив маркера. Если маркер приближается (человек подходит с маркером в руках) --- дрон отлетает назад, сохраняя дистанцию 1 метр. Расстояние берите из \verb|tvec| (нужна калибровка камеры).
\expected{при медленном (до 0,3~м/с) приближении маркера расстояние держится в пределах $1 \pm 0.2$~м; если маркер пропал из кадра, дрон останавливается, а не продолжает движение.}
Опора: \verb|aruco_examples/12_follow_aruco.py|, \verb|aruco_examples/11_aruco_basics_example2_coordinates.py|.
\end{taskbox}

\begin{taskbox}[(К) Инвентаризация]
\emph{Симулятор, затем дрон.} Облетите заданный список точек осмотра, в каждой прочитайте ID в свежем кадре и составьте отчёт <<ID 12 найден в точке (0; 8) в 12:00:05>> (время --- \verb|time.strftime("%H:%M:%S")|). В симуляторе используйте промышленный коридор миссии №\,105 (метки 12 и 13) или карту арены миссии №\,104; на дроне разложите в комнате 5 маркеров с ID от 0 до 4 и допишите запись отчёта в файл.
\expected{в отчёте каждый ID встречается один раз и привязан к точке, в которой он действительно прочитан; ID, которых нет на месте, отчёт называет «не найден». Порядок строк совпадает с порядком облёта.}
\mission{Python}{104--105}
\end{taskbox}

\begin{taskbox}[(К) Посадка на платформу]
\emph{Симулятор, затем дрон.} Дрон должен найти маркер на полу, снизиться, зависнуть точно над центром и совершить посадку.
\expected{посадка не дальше 0,15~м от центра метки; снижение начинается только после трёх кадров подряд с меткой в допуске; при потере метки скорость обнуляется.}
\mission{Python}{98--99}
Опора: \verb|advanced_navigation/staged_landing_sim.py|, \verb|aruco_examples/11_aruco_basics_example3_flight.py|.
\end{taskbox}

\subsection{Блок 2: Телеметрия}

\begin{taskbox}[(Б) Черный ящик]
\emph{Симулятор или дрон.} 10 раз в секунду во время полёта по квадрату печатайте строку CSV: время от старта, $x$, $y$, $z$ (\verb|get_local_position_lps|), курс (\verb|get_yaw|), дальномер (\verb|get_dist_sensor_data|) и напряжение батареи. На дроне пишите строки в файл.
\expected{первая строка --- заголовок, дальше около 10 строк на каждую секунду полёта; если датчик вернул \texttt{None}, в ячейке пусто, а не 0. Таблица открывается в электронной таблице, и по столбцам $x$, $y$ виден квадрат.}
\mission{Python}{64--66, 72}
Опора: \verb|get_telemetry/GET_local_position.py|, \verb|get_telemetry/GET_battery_status.py|.
\end{taskbox}

\begin{taskbox}[(Б) Сигнализация]
\emph{Симулятор или дрон.} Если напряжение батареи падает ниже 3,5~В на банку (для двухбаночного аккумулятора --- 7,0~В, как в задаче «Умная батарея»), дрон должен включить красную подсветку и начать посадку. Логику проверьте функцией \verb|check_battery(voltage)|, вызвав её с 7,9 и 6,9~В до полёта.
\expected{для 7,9~В функция возвращает «норма», для 6,9~В --- «садимся»; в полёте при срабатывании подсветка красная \emph{до} команды \texttt{land()}, и после посадки моторы выключены.}
Опора: \verb|get_telemetry/GET_battery_status.py|, \verb|03_flight_control/LED_hight_change.py|.
\end{taskbox}

\begin{taskbox}[(П) Высотомер]
\emph{Симулятор.} Взлетите на $Z = 9$ над свободной улицей спального района (сцена миссии №\,103) и пролетите над крышей дома в точке $(9; -7)$. Печатайте раз в 0,2~с высоту LPS ($z$) и дальномер. В каких местах они отличаются? На дроне сравните LPS-высоту и дальномер над столом.
\expected{над улицей значения почти равны, над крышей дальномер показывает меньше на высоту дома (больше 5~м), а $z$ не меняется. Вывод: дальномер измеряет расстояние до поверхности под дроном, а не высоту над стартом.}
\mission{Python}{42, 103}
Опора: \verb|get_telemetry/GET_dist_sensor_data.py|.
\end{taskbox}

\begin{taskbox}[(П) Компас]
\emph{Симулятор или дрон.} Напишите программу, которая разворачивает дрон по курсу 0°, 90°, 180° и 270° (в радианах: \verb|math.radians|), не меняя положения, и после каждого поворота печатает \verb|get_yaw()|. Курс отсчитывается в локальной системе, а не от севера.
\expected{четыре строки, в которых измеренный курс отличается от заданного меньше чем на 0,05~рад (с учётом того, что $3\pi/2$ и $-\pi/2$ --- один и тот же курс); положение по $x$, $y$ за время разворотов меняется не больше чем на 0,1~м.}
\mission{Python}{45}
\end{taskbox}

\begin{taskbox}[(К) Виртуальная стена]
\emph{Симулятор, затем дрон.} Задайте границы зоны полета ($X$ и $Y$ от $-2$ до $2$~м). Двигайте дрон командой \verb|set_manual_speed| со скоростью 0,3~м/с наружу; если телеметрия показывает выход за границы, дрон должен остановиться и вернуться в центр.
\expected{дрон останавливается не дальше 0,2~м за границей и возвращается в $(0; 0)$; программа печатает причину остановки («граница X»), а скорость ограничена 0,3~м/с при любой ошибке положения.}
\mission{Python}{68--71}
Опора: \verb|mission_examples/set_manual_speed.py|.
\end{taskbox}

\subsection{Блок 3: Миссии и полет по точкам}

\begin{taskbox}[(Б) Квадрат]
\emph{Симулятор или дрон.} Пролетите по точкам $(0,0) \to (1,0) \to (1,1) \to (0,1) \to (0,0)$ на высоте 1 метр, дожидаясь каждой точки через \verb|point_reached()|.
\expected{в журнале четыре прибытия в нужном порядке, затем посадка; следующая команда отправляется только после прибытия в предыдущую точку.}
\mission{Python}{46--49}
Опора: \verb|03_flight_control/10_02_square.py|.
\end{taskbox}

\begin{taskbox}[(Б) Лестница]
\emph{Симулятор или дрон.} Пролетите вперед на 1 метр, набрав 0.5 метра высоты. Повторите 3 раза (как по ступенькам). Точки вычисляйте в цикле, а не перечисляйте вручную.
\expected{точки $(1; 0; 1.5)$, $(2; 0; 2.0)$, $(3; 0; 2.5)$ при старте с высоты 1~м; финальная высота 2,5~м по LPS с ошибкой меньше 0,1~м.}
\mission{Python}{50}
Опора: \verb|03_flight_control/10_01_forward_back.py|.
\end{taskbox}

\begin{taskbox}[(П) Спираль]
\emph{Симулятор или дрон.} Напишите генератор точек для полета по спирали (расширяющийся круг с набором высоты): радиус растёт от 0,5 до 1,5~м, высота --- от 1 до 2~м за два витка, 12 точек на виток.
\expected{24 точки; расстояние от каждой точки до оси $(0; 0)$ растёт, высота растёт монотонно; дрон проходит их по порядку и садится.}
Опора: \verb|03_flight_control/10_03_circle_math.py|.
\end{taskbox}

\begin{taskbox}[(П) Патруль]
\emph{Симулятор.} Дрон летает между двумя точками $(0; 0; 3)$ и $(0; 8; 3)$, пока в кадре нижней камеры не появится метка 12, но не больше пяти проходов. На дроне вместо метки можно взять препятствие по дальномеру.
\expected{после первого кадра с меткой 12 патруль прекращается и дрон садится; если метки нет, программа печатает «Патруль завершён» после пяти проходов. Бесконечного цикла нет.}
\mission{Python}{35, 102}
\end{taskbox}

\begin{taskbox}[(К) Сканирование поля]
\emph{Симулятор.} Реализуйте полет <<змейкой>> над зоной $8 \times 8$~м на высоте 3~м так, чтобы кадры нижней камеры перекрывали всю зону (с $Z = 3$ камера видит по $Y$ около $\pm 2$~м). В каждой точке снимайте кадр и печатайте все прочитанные ID. На дроне с той же логикой получают фотоплан.
\expected{маршрут из параллельных проходов с шагом не больше 4~м; итоговый список ID без повторов; программа печатает число точек съёмки и время облёта.}
\mission{Python}{101}
Опора: \verb|mission_examples/mission_with_frames.py|.
\end{taskbox}

\subsection{Блок 4: Групповой полет}

\begin{taskbox}[(Б) Синхронный старт]
\emph{Дрон или симулятор с двумя дронами.} Два дрона взлетают одновременно по команде с ПК, висят 5 секунд и садятся. В симуляторе добавьте второй дрон в список дронов и возьмите его адрес из настроек подключения.
\expected{разница между командами взлёта двух дронов меньше 0,5~с; оба дрона садятся, и программа закрывает оба соединения.}
Опора: \verb|group_flight_examples/swarm_OPT.py|.
\end{taskbox}

\begin{taskbox}[(П) Карусель]
\emph{Дрон или симулятор с двумя дронами.} Два дрона висят напротив друг друга и меняются местами, пролетая по дуге (чтобы не столкнуться).
\expected{в течение всего манёвра расстояние между дронами по телеметрии не меньше 1~м; каждый дрон заканчивает в стартовой точке другого.}
Опора: \verb|group_flight_examples/swarm_circle_IRNAV.py|.
\end{taskbox}

\begin{taskbox}[(П) Ведущий-ведомый]
\emph{Дрон или симулятор с двумя дронами.} Первый дрон летит по программе (квадрат). Второй дрон получает телеметрию первого и повторяет его движения со смещением (на 1 метр сзади).
\expected{после каждой вершины квадрата ведомый находится в 1~м позади ведущего с ошибкой меньше 0,3~м; если телеметрия ведущего не пришла (\texttt{None}), ведомый висит на месте.}
Опора: \verb|group_flight_examples/swarm_OPT.py|.
\end{taskbox}

\begin{taskbox}[(К) Шоу дронов]
\emph{Дрон или симулятор.} Составьте хореографию для 4 дронов под музыку (используйте мигание светодиодами в такт).
\expected{сценарий записан таблицей «время --- точка --- цвет» для каждого дрона; ни в один момент два дрона не ближе 1~м друг к другу; смена цвета попадает в такт с ошибкой меньше 0,2~с.}
Опора: \verb|group_flight_examples/swarm_OPT.py|, глава о групповом полёте.
\end{taskbox}

\subsection{Блок 5: Трекинг объектов}

Распознавание людей в симуляторе не моделируется: задания блока выполняются с настоящим дроном или с веб-камерой компьютера.

\begin{taskbox}[(Б) Приветствие]
\emph{Дрон.} Если дрон видит лицо человека --- он мигает зеленым и делает <<кивок>> (короткое движение вперёд-назад).
\expected{кивок выполняется один раз на появление лица, а не на каждом кадре; без лица в кадре дрон неподвижен.}
Опора: \verb|PioneerHumanTracking/human-tracking.py|.
\end{taskbox}

\begin{taskbox}[(П) Следи за цветом]
\emph{Дрон.} Возьмите яркий оранжевый мяч. Дрон должен поворачиваться (yaw) так, чтобы мяч всегда был в центре кадра.
\expected{при медленном перемещении мяча центр пятна остаётся в пределах 10\% ширины кадра от центра; если мяч пропал, вращение прекращается.}
Опора: \verb|camera_examples/camera_stream.py|.
\end{taskbox}

\begin{taskbox}[(П) Держи дистанцию]
\emph{Дрон.} Дрон видит лицо. Если лицо занимает слишком много места в кадре (человек близко) --- отлетаем. Слишком мало (человек далеко) --- подлетаем.
\expected{в установившемся режиме высота рамки лица держится в заданном диапазоне; скорость ограничена (не больше 0,3~м/с) и обнуляется в мёртвой зоне.}
Опора: \verb|PioneerHumanTracking/human-tracking.py|.
\end{taskbox}

\begin{taskbox}[(К) Оператор]
\emph{Дрон.} Дрон летит боком параллельно идущему человеку, снимая его профиль (требует сложной координации движений).
\expected{человек всё время в кадре, а боковая скорость дрона следует за его шагом; при потере человека дрон останавливается и висит.}
Опора: \verb|PioneerHumanTracking/human-tracking.py|.
\end{taskbox}

\subsection{Блок 6: Самостоятельные миссии симулятора}

Последние девять миссий панели \ui{Обучение} (раздел «Самостоятельные миссии») собирают всё, что было в части о зрении: маршрут, камеру, детектор, наведение и отчёт. Проверку выполняет сам симулятор по журналу полёта: важны фактические прибытия, свежие кадры и прочитанные ID, а не напечатанное слово «готово».

\begin{taskbox}[Карта из нескольких меток]
\emph{Симулятор.} Свяжите ID с назначением словарём: 10 --- «северо-запад», 11 --- «юго-восток», 12 --- «цель». Напечатайте назначение только тех ID, которые реально прочитаны в кадре, и посадите дрон.
\expected{три строки «ID:назначение» из текущей детекции; ID, которого нет на карте, не печатается.}
\mission{Python}{100}
\end{taskbox}

\begin{taskbox}[Поиск по секторам]
\emph{Симулятор.} Найдите ID 12, проверяя три сектора вдоль $Y$ с шагом 4~м начиная с $Y = 0$ ($X = 0$, высота 3~м); после находки прекратите поиск, вернитесь и сядьте.
\expected{проверены два пустых сектора и сектор с целью; после находки секторов больше нет; если цели нет нигде --- «Поиск истёк».}
\mission{Python}{101}
\end{taskbox}

\begin{taskbox}[Лесная экспедиция]
\emph{Симулятор.} В лесном полигоне по свободному центральному коридору (линия $X = 0$) долетите до метки 12 в 8~м от старта, прочитайте её ID, вернитесь и сядьте.
\expected{точка осмотра достигнута, ID 12 прочитан из кадра в этой точке, обратный путь и посадка завершены.}
\mission{Python}{102}
\end{taskbox}

\begin{taskbox}[Осмотр крыш]
\emph{Симулятор.} В спальном районе осмотрите крыши домов в $(9; -7)$ и $(-9; -7)$ с высоты $Z = 9$, прочитайте ID 102 и 101 и вернитесь на старт по верхнему маршруту.
\expected{оба крышных ID прочитаны над своими домами; высоту набирают до горизонтального перелёта.}
\mission{Python}{103}
\end{taskbox}

\begin{taskbox}[Контрольные точки арены]
\emph{Симулятор.} На Геоскан Арене подлетите сверху к карте ArUco в $(-9; 10.5)$ и осмотрите её с высоты 4~м из двух точек: прямо над ней и со сдвигом 0,5~м по $+X$. Напечатайте ID, прочитанные в каждой точке.
\expected{обе точки достигнуты; в них прочитаны ID 10 и 11; возврат через $Z = 6$ и посадка.}
\mission{Python}{104}
\end{taskbox}

\begin{taskbox}[Инспекция промышленной зоны]
\emph{Симулятор.} Осмотрите два участка промышленного коридора: метка 12 при $Y = 8$, метка 13 --- на 8~м дальше по той же линии $X = 0$. Для каждого участка напечатайте прочитанный ID, затем вернитесь и сядьте.
\expected{12 и 13 прочитаны в своих точках и в правильном порядке; к следующему участку дрон летит только после подтверждения текущего.}
\mission{Python}{105}
\end{taskbox}

\begin{taskbox}[Неожиданная потеря цели]
\emph{Симулятор.} Получите цель, намеренно потеряйте её отлётом к $Y = 5$ и восстановите не более чем за три попытки (на второй --- возврат к старту); после обнаружения прекратите поиск и сядьте.
\expected{в журнале есть цель до потери, пустой результат, нулевая скорость и новое обнаружение; без возврата программа сообщает об исчерпании попыток.}
\mission{Python}{106}
\end{taskbox}

\begin{taskbox}[Незнакомая карта]
\emph{Симулятор.} Наведитесь и сядьте на ID 12 на смешанной карте тем же алгоритмом, что в миссиях №\,96--99, не записывая координаты метки в код.
\expected{посадка над меткой 12; точка посадки отличается от привычной $(0.6; 0.3)$, потому что метка на этой карте стоит в другом месте.}
\mission{Python}{107}
Опора: \verb|advanced_navigation/staged_landing_sim.py|.
\end{taskbox}

\begin{taskbox}[Итоговая экспедиция]
\emph{Симулятор.} Соберите полную миссию: поиск ID, намеренная потеря и восстановление, наведение и посадка. Первый запуск --- с ID 12; второй --- с ID 249, которого на карте нет.
\expected{первый запуск заканчивается посадкой на метку 12 после восстановления; второй --- сообщением «Поиск истёк», возвратом и посадкой без ложного успеха. Обе ветки проверяет кнопка \ui{Проверить полёт}.}
\mission{Python}{108}
\end{taskbox}

\subsection{Критерии оценки выполнения}
Задание считается выполненным, если:
\begin{enumerate}
    \item Код запускается без синтаксических ошибок.
    \item Выполнено условие из строки «Что должно получиться»; в симуляторе --- на двух разных seed полигона или после сброса сцены, на дроне --- в реальных условиях не менее 3 раз подряд.
    \item Предусмотрены неуспешные исходы: метка не найдена, кадр не получен (\texttt{None}), точка не достигнута за отведённое время. Программа сообщает причину и сажает дрон, а не висит бесконечно.
    \item Код содержит комментарии, объясняющие логику работы.
    \item Соблюдены меры безопасности (дрон не вылетает за пределы зоны, экстренная посадка работает).
\end{enumerate}
